\chapter{Complex Challenges \& Public Engagement}

\openepigraph{%
	A ruined planet cannot sustain human lives in good health. A healthy planet and healthy people are two sides of the same coin.
}{Dr. Margaret Chan, Executive Director, WHO}

\marginnote{
	\newthought{Key Message:} The impact of climate heating on local populations is a truly wicked problem that public health practitioners spend very little time thinking about. It's too hard. It isn't a health system priority; it is tomorrow's problem. We need to make it easier to understand and incorporate into all local population planning. Systems thinking might provides some useful tools.
}


\section{Executive Summary}
The aim of this document is to bring together all threads of this project that set out to explore how public health could develop policy in the area of potential health impacts from climate change on the older population of West Dorset. It became obvious quite early in the project that this was not going to be an easy task. The difficulty arises firstly in developing a meaningful engagement on the topic. Getting attention for a problem that is in the future, and not thought to be more significant in impact than many current population health issues was always going to be a challenge.

We have approached this challenge by adopting a `systems thinking' approach to our engagement with policy makers and `futures studies' approach in our engagement with the community more widely. In truth, these approaches are complementary and underpin both engagement processes.

\begin{description}
\item[Systems Thinking]An approach to developing deeper understanding of complexity by examining problems and challenges within the context of the `system' in which they exist.
\item[Futures Studies] An approach to the development of strategy that seeks to understand an issue or challenge in the context of a more comprehensive picture of a future state. Incorporates a fundamental belief that at any point in time there are many possible futures -- strategically one should work towards a preferred future.
\end{description}

\section{Change of Project Emphasis}
The project that began as a climate change impacts on the health of older persons in West Dorset, ended as a project entitled, \emph{Complex Challenges \& Public Engagement}, reflecting the broader issues we encountered during the running of this project.

Our change in emphasis was not an abandonment of the initial research focus that sought to develop a better understanding of what impacts older people in West Dorset would face in thirty years. Indeed, it was our desire to make a meaningful contribution in this area that we realised that our initially more traditional research approach had two shortcomings:
\begin{enumerate}
\item Scientific uncertainties regarding the estimated outcomes 30 years in the future are such that we would not succeed in delivering new information, but rather our approach would include the same lists of possible outcomes and generalised mortality outcomes; and
\item Consequently, the initial difficulties we experienced in engaging stakeholders would be compounded, because we had nothing particularly compelling to engage with.
\end{enumerate}

\section{Project Outputs}
Thus far this project has presented a number of workshops, a conference poster and delivered a systems model and story, various futuring deliverables including a Futures Story and this report.

It is anticipated that these materials will be used further in 2016 to continue exploring the public health challenge faced by climate change in Dorset and we will also seek to publish some of this work.

Furthermore, a longer term outcome will require a longer effort as we attempt to apply the methods explored here to other public health challenges over the coming years.

\section{Project Outcomes}

Ultimately, Public Health Dorset role is to protect and improve the health of its population -- we hope that methods adopted during this project will contribute towards the success of that role. However, a more immediate outcome has been the exposure of many staff to systems thinking and futuring methods, which now extends their potential reach in developing shared understanding of particularly complex public health challenges.

\section{Report Organisation}

Part I of this report focusses on the `systems thinking', `systems modelling' and `systems story telling' methods that were used in an effort to develop a shared understanding of this public health challenge.

Part II reports on the use of future methods: `horizon scanning', `scenario building' and `futures story telling'.

Part III of this report serves as a repository for presentations and stakeholders events, recording the evolution of the group's thinking over the two year life of the project.